\documentclass[11pt]{article}
\usepackage[letterpaper,margin=1in]{geometry}
\usepackage[T1]{fontenc}
\usepackage{lmodern,amsmath,amssymb,amsthm,booktabs,longtable,array,hyperref}
\hypersetup{colorlinks=true,linkcolor=blue,urlcolor=blue}
\newtheorem{theorem}{Theorem}[section]
\newtheorem{lemma}[theorem]{Lemma}
\newtheorem{proposition}[theorem]{Proposition}
\newtheorem{corollary}[theorem]{Corollary}
\theoremstyle{definition}\newtheorem{definition}[theorem]{Definition}
\theoremstyle{remark}\newtheorem{remark}[theorem]{Remark}
\newcommand{\Z}{\mathbb Z}\newcommand{\N}{\mathbb N}
\newcommand{\Con}{\operatorname{Con}}
\setlength{\emergencystretch}{3em}
\title{Unsquared factors and linear selectors for Diophantine consistency}
\author{Research prepared for Vladimir Reshetnikov}
\date{25 October 2026}
\begin{document}
\maketitle
\begin{abstract}
We give new explicit polynomial equations whose unsolvability in their
stated domains is equivalent in ZFC to syntactic consistency of ZFC.
The main natural-number construction evaluates an unsquared product of
proof-checking factors instead of separately squaring every factor.
A canonical CRT choice makes all newly exposed factors nonnegative,
while preserving the semantic anchor and the exact source polynomial.
The unrestricted-integer construction combines a signed radix with a
linear binomial selector radius, removing redundant arithmetic witnesses.
For ten positive variables, an inverse Lucas correction reuses existing
registers but requires a new projection-first divisibility argument.
A separate root-certificate circuit is shortened by exact identities
and an integer coordinate shear.
The release supplies literal equations, full fixed numerals, ordinary
proofs, inverse maps, reproducible builders and independent checks.
All statements retain their inherited calibration boundaries. No
consistency decision, global minimum or new end-to-end kernel certificate
is claimed.
\end{abstract}
\tableofcontents
\section{Objective and comparison scope}
For each fixed polynomial $P\in\Z[x_1,\ldots,x_n]$ and specified domain $D$,
the target is
\[
 \mathrm{ZFC}\vdash
 \bigl(\neg\exists\boldsymbol{x}\in D^n\;P(\boldsymbol{x})=0\bigr)
 \longleftrightarrow\Con(\mathrm{ZFC}).
\]
We prove solvability equivalent to existence of a finite syntactic
refutation. We do not prove either consistency or unsolvability separately.
The domains $\Z$, $\N=\{0,1,2,\ldots\}$ and $\Z_{>0}$ are part of the
individual theorems and are not interchangeable complexity measures.

The comparison snapshot is
\texttt{7a8da414376a80b90ecf799f37ceb26fb9ac674e} of
\href{https://github.com/VladimirReshetnikov/ProveIt}{ProveIt}.
It includes the new factorial-free, shared factorial-error, nearest-factorial,
normalized-residue, digit-recovery and Lucas-correction incoming packages,
as well as our preceding cofinal-witness report. We compare with their
actual stronger equations, not only with our own older results. Incoming
placement, repository adoption, ordinary proof review and fresh finite
replay are recorded separately in the baseline audit.

The measurements are quantified coordinates, degree, quadratic residual
count, fully collected support, maximum absolute coefficient and live
binary evaluator operations. Addition, subtraction and multiplication
cost one each, including squares and scalar products; register reuse is
free. Fixed coefficient ports are not unknowns, but their entire coherent
numerical tuples and finite definitions are part of every specification.
Circuit degree and operation count do not measure expanded coefficient
height or bit complexity.

\subsection{Selected explicit equations}
\begin{center}\small
\begin{tabular}{lrrrrr}\toprule
Domain and profile & Variables & Degree & Rows & Terms & Height\\\midrule
$\N$, dimension &206&4&158&8210&130686376\\
$\N$, balanced &207&4&159&5278&69238087\\
$\N$, sparse &208&4&160&5198&65866712\\\midrule
$\Z$, dimension &239&4&175&12658&6144\\
$\Z$, balanced &240&4&176&5232&4096\\
$\Z$, sparse &240&4&176&4927&7596\\\bottomrule
\end{tabular}
\end{center}
The natural 206-variable equation improves the matched preceding
219-variable/171-row/8267-term/height-130686428 equation in all four
varying metrics. The sparse natural point similarly improves the
preceding 221-variable/173-row/5253-term/height-65867916 equation.
The middle point provides a separate dimension/support/height tradeoff.

The new 239-variable integer quartic improves the incoming shared-error
240-variable/177-row/21482-term/height-6144 point in dimension, rows and
support at the same height. The sparse integer point improves the
incoming factorial-free 242-variable/178-row/5018-term/height-7968
comparison in every varying metric. The balanced point trades support
for a smaller coefficient height; these comparisons are not global
variable-count records.

\begin{center}
\begin{tabular}{lrrr}\toprule
Domain and circuit & Variables & Exact degree & Binary operations\\\midrule
Strictly positive, inverse correction &10&304&85\\
Unrestricted, completed-square carrier &13&152&119\\\bottomrule
\end{tabular}
\end{center}
The positive circuit improves the incoming degree-304/87-operation
Lucas-correction form at unchanged degree and domain. Its
84-operation/degree-410 square-weight companion is a dominated control.
The 13-integer circuit improves the matched degree-152/126-operation
incoming form. It remains a tradeoff against our earlier degree-156/
117-operation point. Its coordinate shear is an exact integer bijection,
not an identity in the original unshifted coordinate names.

Both numerical specializations of the positive circuit are supplied:
the original radix67 tuple and the complete previously reviewed
incoming complemented-label radix47 tuple. The latter has 207214 raw
coefficient bytes instead of 295388; its encoding is credited prior
work, not a new encoding theorem in this report. Counts and degrees
are unchanged, while the specialized integer polynomials differ.

\section{Canonical CRT signs and common-square extraction}
\label{sec:round20-natural}

Every coordinate in this section independently ranges over
\(\mathbb N=\{0,1,2,\ldots\}\). The selected ordinary quartics are
\(P(\mathbf z)=\sum_j R_j(\mathbf z)^2\), with every quadratic residual
and every expanded integer coefficient supplied literally in the packet.
\begin{center}
\begin{tabular}{lrrrr}
\hline
Endpoint & Variables & Rows & Quartic terms & Height\\
\hline
Previous dimension control &219&171&8,267&130,686,428\\
Previous sparse control &221&173&5,253&65,867,916\\
New \texttt{natural206} &206&158&8,210&130,686,376\\
New \texttt{balanced207} &207&159&5,278&69,238,087\\
New \texttt{sparse208} &208&160&5,198&65,866,712\\
\hline
\end{tabular}
\end{center}
The new dimension and sparse equations improve their corresponding controls
in all four displayed metrics; the 207-variable equation is a middle tradeoff.
These are measured constructions, not claims of global optimality.
Relative to the retained ordinary calibration from C222 through ZF to ZFC,
each has a natural
zero exactly when a finite syntactic ZFC refutation exists. Hence ZFC proves
its no-natural-zero sentence equivalent to \(\operatorname{Con}(\mathrm{ZFC})\).
This is neither a consistency proof nor a newly kernel-checked end-to-end
logical derivation. Factorials, Pell sequences and powers below belong to
the ordinary proof, not to the syntax of the fixed released quartics.

\subsection{An exact factorization of the unchanged local source}

Use the old source's representative-minus-one coordinates \(a,b,c,d,S,S'\).
Write
\begin{align*}
 \delta_0&=S-S',& \delta_i&=S-S'-m_i\quad(1\le i\le8),\\
 \delta_{11}&=(a-c)^2(b-c)^2-(d+2)^2,
 &\delta_{12}&=S-S'-\texttt{mul.358},\\
 \delta_{13}&=\texttt{mul.363}-\texttt{mul.369},
 &\delta_{14}&=\texttt{mul.372}-\texttt{mul.373}.
\end{align*}
Here \((m_1,\ldots,m_8)\) are the already existing gates
\[
 (\texttt{mul.339},\texttt{mul.341},\texttt{mul.343},\texttt{mul.345},
 \texttt{mul.347},\texttt{mul.349},\texttt{mul.351},\texttt{mul.353}).
\]
The old axiom/rule disjunction is the literal nonnegative polynomial
\[
 R=\left(\prod_{i=0}^{8}\delta_i^2\right)
       (\delta_{11}^2+\delta_{12}^2)\delta_{13}^2\delta_{14}^2.
 \tag{C1}
\]
The complete local source adds the two unchanged beta-difference squares,
\(R+\beta_1^2+\beta_2^2\). No change to the logical calculus is being made.

Define eleven oriented affine factors and two useful expressions by
\[
 g_i=-\delta_i\ (0\le i\le8),\qquad
 g_9=-\delta_{13},\qquad g_{10}=\delta_{14},\qquad
 Q=\prod_{i=0}^{10}g_i,\quad E=(a-c)(b-c),\quad D=d+2.
\]
Ten quadratic product rows construct \(Q\) successively, starting with
\(g_0g_1\); affine factors require no new coordinate. One row constructs
\(E\). Three more quadratic rows construct
\[
 Z_0=Q(E-D),\qquad Z=Z_0(E+D),\qquad T=-Q\delta_{12}.
 \tag{C2}
\]
Replace the old evaluation input \(R\) by \(Z^2+T^2\), retaining the full
source-evaluation row and its natural quotient. The integer identities
\[
 Q^2=\left(\prod_{i=0}^{8}\delta_i^2\right)\delta_{13}^2\delta_{14}^2,
 \qquad (E-D)(E+D)=E^2-D^2=\delta_{11}
\]
give \(Z^2+T^2=R\), without any sign assumption.
Thus \emph{the entire local source polynomial is unchanged}, as are its
majorant, prime-power decoding and logical interpretation.

The old fifteen square gates, two inequality-product gates and eleven
disjunction-product gates become fourteen product gates: a saving of
fourteen coordinates and fourteen rows before contraction. The new raw
source has \(310\) natural coordinates and \(262\) quadratic rows. All
thirteen nonanchor Pell seeds, eleven digit requests, six exact powers,
source endpoints, semantic frame and derived label polynomial are retained.

\subsection{Pointwise soundness and the genuine natural-domain issue}

From any natural zero of the new source, restore each old square as its
literal square and restore the old nonnegative product chain in order.
These values are natural regardless of the signs of the affine factors.
The identity above restores the old evaluation row; every other old row
was retained verbatim. We have therefore obtained a natural zero of the
complete old source, to which its established soundness proof applies.

Conversely, an arbitrary old tuple need not make the new partial products
nonnegative. Consequently the transformation is \textbf{not claimed to be
an orthant automorphism or a pointwise extension of all old witnesses}.
The required completeness direction starts from a finite refutation and
uses the freedom in choosing its canonical CRT representatives.

\subsection{Cofinal sign choices preserve every CRT condition}

First perform the parent's even descending construction, before choosing
the encoded column representatives:
\[
 t=2t_0+4,\quad F=t!,\quad n=F^2-1,\quad V=n^t,\quad y=V-1.
\]
Choose \(t\) larger than the finite proof-tree data and all its local labels.
There are at most \(t\) pending nontruth vertices, each with weight at most
\(1+tF<n\), so every pending state is at most \((1+tF)^t<V\), hence at
most \(y\). Since \(t\) is even, \(F\mid y\). The beta moduli
\(1+(i+1)y\), \(0\le i\le t\), are pairwise coprime: a common prime
would divide a positive index difference at most \(t\), and thus \(F\)
and \(y\), contradicting its division of either modulus. Beta CRT gives
\(x\) with the prescribed pending states and endpoint residues; enlarge
\(x\) within its class as required. Define
\[
 B=\binom nt,\qquad N=B^x.
\]
In particular, \(N\) has been fixed before any subsequent lift.

The inherited valuation construction uses positive representatives modulo
\(N^2\). Adding an arbitrary multiple of \(N^2\) to any minus-one
coordinate preserves its required residue, positivity and precision; there
is no parity or square restriction on these natural coordinates.
Choose the representatives for \(c,d,S,u,v\) and the encoded row index
first. Next choose \(b>c\) in its prescribed class and then \(S'>S\)
in its prescribed class. Finally choose \(a\) in its class so large that
\begin{gather*}
 a>c+d+3,\qquad
 a\ge\text{every previously fixed nonanchor index and coordinate},\\
 (S+1)(1+F(a+1))>(S'+1)(1+F(b+1)).
 \tag{C3}
\end{gather*}
All these are lower bounds on \(a\); a full tail of its arithmetic
progression satisfies them. There is no dependence of \(N\) on these choices.

We now verify every new factor's sign. The factors \(g_0=S'-S\) and
\(g_i=S'-S+m_i\), \(1\le i\le8\), are positive, as is
\(-\delta_{12}=S'-S+\texttt{mul.358}\). To check the two remaining
factors, set \(M=\texttt{mul.36}\). The literal inherited rows give
\begin{align*}
 \texttt{mul.363}&=(S'+1)(1+F(b+1)),\\
 \texttt{mul.368}&=(S+1)(1+F(a+1)),\\
 \texttt{mul.369}&=(1+F(4+4b+2M))\texttt{mul.368}.
\end{align*}
Thus (C3) implies \(g_9=\texttt{mul.369}-\texttt{mul.363}>0\).
Also
\begin{align*}
 \texttt{mul.372}&=(S'+1)(1+F(6+4b+2M)),\\
 \texttt{mul.373}&=(S+1)(1+F(b+1)).
\end{align*}
Their inner weight difference is \(F(5+3b+2M)>0\); combined with
\(S'>S\), this gives \(g_{10}>0\). Finally \(b-c\ge1\) and
\(a-c>d+3\) imply \(E>D>0\). Every partial product, \(Z_0,Z,T\)
in (C2) is therefore natural.

The ordinary local residues are unchanged by these lifts, so the inherited
precision argument still makes the nonnegative source a multiple of \(N\).
Its quotient is natural. This also holds for its exactly equal new expression.

\subsection{The anchor is chosen after the signs, without a cycle}

There is an important constraint beyond CRT residues: the semantic Pell
anchor must dominate every requested seed index. It is not enough merely
to enlarge the self-sized radix. Here the actual source anchor is unchanged:
\[
 J=7+4b+2M+\texttt{add.402},\qquad
 M=(a+b)^2+5(a+b)+6.
\]
All remaining terms are nonnegative, so
\(J\ge2(a+1)^2\). Condition (C3) makes every nonanchor seed index at
most \(a+1\), including the index \(a+1\) itself. Hence
\[
 J\ge2\max_i k_i,
\]
exactly the inherited sufficient condition for natural completeness of the
one-row error-square cap. No lift of \(S\) or \(S'\) occurs afterward;
neither is required to exceed \(J\) or a derived syntax weight.

Only now define the positive sum \(S_{\rm base}\) of the independent
representatives, with its fixed offset, and choose
\(q=(NS_{\rm base})^x\). Every representative is below \(S_{\rm base}\),
whatever the size of its lift. The unchanged constructor ledger gives
\(J\le S_{\rm base}^{\ell}\), \(\ell=2^{22}\). The inherited even-horizon
bootstrap gives \(y>2^{36}\), \(x>2^{38}\), and \(\ell(t+1)<x\).
Thus all required nonfactorial low coefficients are at most
\(S_{\rm base}^x\), and factorial low coefficients are at most \(V<N\).
The four pooled digit groups, of sizes \(2,5,2,2\), remain carry-free since
\[
 5S_{\rm base}^x+N<N^xS_{\rm base}^x=q.
\]
Their low-remainder sums are at most one less than the respective denominator,
so all pooled slacks are natural.

Compute the remaining denominator powers and frame \(W\) next; choose its
lower Pell witness, a sufficiently large \(A\) in the common-root class,
and then all principal, positive-center auxiliary, comparison and seed Pell
witnesses. The anchor inequality supplies the natural merged-error quotient.
None of these later choices changes the already fixed indices or CRT lifts.
This completes the acyclic construction of a new natural zero.

\subsection{Independent finite certification and artifacts}

The source checker independently identifies all 28 removed gate definitions,
all 14 new gate definitions, the evaluation-row replacement and all 247
otherwise unchanged rows. It verifies the factorization compositionally by
exact integer identities, not by random evaluations. Corruption controls
include deleting the source-evaluation row, changing a factor, widening the
domain to integers and modifying an unrelated row.

The subsequent monic eliminations are separately and pointwise certified:
every eliminated coordinate is an integer polynomial equal to a positive
rational sum of affine squares plus nonnegative monomials. Independent
replay proves their natural restorability and every substitution. They yield
\(206/158\), \(207/159\) and \(208/160\) controls. Finally 25, 30 and 33 invertible integer
\emph{row} shears produce the selected endpoints in the table; coordinates
and their natural domains do not change. Exact inverse operations, every
expanded coefficient, refreshed metrics and printed residuals parsed back
from TeX are verified both normally and with Python assertions disabled.

The full proof, complete literal quartics, pre-basis controls, certificates,
offline rebuild and frozen ordinary arithmetic/logical closure are included
under \texttt{natural/}. The new theorem uses canonical sign choices only
for completeness; soundness and the later certified contractions remain
pointwise. This distinction is essential to the reduction.

\section{A linear binomial radius and a signed radix in one residue class}
\label{sec:r20-direct}

All equations in this section are ordinary quartics over independently
unrestricted integers.  They use the numerical C222 proof calculus directly,
not a universal-polynomial machine-index substitution.  Relative to the
incoming factorial-free construction, two semantic edits eliminate two
coordinates and two residuals:
\[
 R=F^2\quad\longrightarrow\quad R=F,
 \qquad
 q=(2z+1)^2\quad\longrightarrow\quad q=4z+1,
 \qquad F=[2(t+1)]^t,\quad t=16+h^2.
\]
Here the variable power describes the proved meaning of an unchanged
quadratic Pell certificate; it is not an operation in the final quartic.
The first edit uses a classical lcm estimate.  The second allows an
initially negative radix and then proves its positivity in the necessary
order.  Exact residual-basis changes subsequently improve the collected
coefficients without changing the zero set.

Each final equation has the form
\[
 P(\boldsymbol x)=\sum_{i=0}^{m-1}R_i(\boldsymbol x)^2=0,
 \qquad\boldsymbol x\in\mathbb Z^n,
\]
with all $R_i$ of degree at most two and $P$ of exact degree four.  Every
residual and every collected coefficient is supplied, together with
coordinate and inverse row certificates.  The accompanying endpoint table
records the measured tradeoffs.  ``Height'' always means the largest
absolute coefficient after complete collection of $P$.
\begin{center}
\begin{tabular}{lrrrr}
\toprule
Profile & Variables & Rows & Monomials & Height\\
\midrule
Dimension & 239 & 175 & 12,658 & 6,144\\
Balanced & 240 & 176 & 5,232 & 4,096\\
Sparse & 240 & 176 & 4,927 & 7,596\\
\bottomrule
\end{tabular}
\end{center}
The 239-coordinate profile improves variables, residuals and collected
support over the incoming Shared Factorial Error 240-coordinate endpoint,
while retaining its height 6,144.  All three new profiles have maximum
residual coefficient 64.  These are achieved tradeoffs, not global minima.

\subsection{A classical lcm estimate makes the radius linear}
Write $\Lambda_n=\operatorname{lcm}(1,\ldots,n)$.  We recall the classical
bound $\Lambda_n\le4^n$, including a short proof so that the improvement
does not depend on an unexplained asymptotic estimate.  For $n\ge2$, set
$k=\lceil n/2\rceil$.  Then
\begin{equation}\label{eq:r20-lcm-recur}
 \Lambda_n\mid\Lambda_k\binom nk.
\end{equation}
Indeed, let $p^a$ be the largest power of $p$ not exceeding $n$.
If $p^a\le k$, it already divides $\Lambda_k$.  Otherwise
$p^{a-1}\le k$, so only one further factor $p$ is needed.
The $p^a$ term of Legendre's valuation formula for $\binom nk$ is one:
$n$ has one multiple of $p^a$, while $k$ and $n-k$ have none.
All other terms of that formula are nonnegative.

For $n=2m$, the middle binomial coefficient is at most $4^m$.
For $n=2m+1$, its two equal middle coefficients have sum at most
$2^{2m+1}$, so each is again at most $4^m$.  Complete induction in
\eqref{eq:r20-lcm-recur}, starting with $\Lambda_1=1$, yields
$\Lambda_n\le4^n$ in both cases.  This is a classical theorem, not a
claimed new number-theoretic result; see Chan and Norrish,
\emph{Proof Pearl: Bounding Least Common Multiples with Triangles},
J.~Automated Reasoning 62 (2019), 171--192, Theorem~3,
\href{https://doi.org/10.1007/s10817-017-9438-0}{doi:10.1007/s10817-017-9438-0}.

The unchanged horizon satisfies $t\ge16$.  Consequently
\begin{equation}\label{eq:r20-linear-radius}
 F=[2(t+1)]^t\ge34^t>16^t\ge\Lambda_t^2.
\end{equation}
The authoritative factorial-free predecessor, its
\texttt{report/main.tex}, proves the following two-sided binomial selector
lemma under the hypothesis $R\ge\Lambda_t^2$:
for $B=\binom{R+t}{t}$, every row $i\in\{1,\ldots,t\}$ has a prime power
$p^e$ exactly dividing $B$ with $p^e>t$ and $p^e\mid R+i$; conversely,
every prime power exactly dividing $B$ divides $R+i$ for some row.
The first assertion recovers every local proof row.  The second handles
\emph{all} primes in the reverse CRT construction, including primes not
selected to isolate a row.  Equation~\eqref{eq:r20-linear-radius} therefore
permits $R=F$.  An older predecessor review proves only a factorial-square
threshold; that weaker review is not the theorem used here.

\subsection{What is retained before signs are recovered}
The complete signed squared frame, semantic anchor, pooled Pell orientation
conditions, and all eighteen power slots are retained.  Both coordinates
of every principal Pell pair are oriented at their signed requested
indices.  No representative subset or frame atom is discarded.
The inherited signed-power lemma, using the strict frame bound
$|A|>W^W$, states that at a nonzero integer base $b$ with $|b|\ge2$, a
requested index $k$ must be nonnegative and its output equals $b^k$.
For $b=\pm1$, the output is $1$ or $(-1)^k$ even for signed $k$.
For a negative index at a nonunit base, the same small-error congruence
would force an integer identity $m b^{-k}=1$, which is impossible.
We never apply this lemma to the base zero.

The known positive index $t$ and base $2(t+1)$ first give
$F=[2(t+1)]^t>2^{32}$, independently of any radix sign.  The retained
beta certificate
\[
 v_\beta=(2w+1)^2,\qquad y=Fv_\beta
\]
then yields $y\ge F>0$.  The possible boundary $y=F$ is not silently
excluded; it is addressed after genuine local-row recovery below.

\subsection{The mod-four radix and the acyclic bootstrap}
The replacement $q=4z+1$ excludes both $q=0$ and $q=-1$.
It is stronger than oddness alone in precisely the needed respect:
$q=-1$ would make the numerator base $q+1$ zero.
We now exclude the only remaining unit, $q=1$.
In that case the base $q+1=2$ makes all nine actual CRT representative
exponents nonnegative.  Hence
\[
 H=2+\sum_\nu a_\nu\ge2,
 \qquad C_0\ge0.
\]
The retained raw index equation $B\mu=C_0+R+1$ has positive right-hand
side, so $B\ne0$ and $|BH|\ge2$.
The self-radix request gives $x\ge0$ and $1=(BH)^x$, forcing $x=0$.
But the last beta endpoint is
\[
 x=1+[1+(t+1)y]u_t.
\]
At $x=0$ it would make the integer $1+(t+1)y>1$ divide $-1$.
This contradiction uses neither a binomial interpretation nor a remainder
bound.

For every remaining radix, put $Q=|q|$.  Then $Q\ge3$ and
$|q+1|\ge2$.  The same argument recovers nonnegative representatives,
$H\ge2$, and $B\ne0$.  The derived label is another requested exponent
at $q+1$, so it is nonnegative; the unchanged anchor excludes zero.
The genuine self-radix equality is $q=(BH)^x$ with $x\ge0$.
Since $x=0$ would give $q=1$, we have $x\ge1$.
If the first beta quotient were $u_0\le-1$, its endpoint
\[
 x=F+1+(y+1)u_0
\]
would imply $x\le F-y\le0$.  Thus $u_0\ge0$, $x\ge F+1$, and the
last endpoint has $u_t\ge1$.  It follows that
\begin{equation}\label{eq:r20-q-size}
 x\ge2+(t+1)y\ge2+(t+1)F,\qquad
 Q=|BH|^x\ge2^x,\quad Q\ge4B^2,\quad |B|\le Q/4.
\end{equation}
No sign of $q$ or $B$ has yet been assumed.

All powers at $q$ and $q+1$ now have their genuine meanings at their
nonnegative indices.  In particular $M=q^x\ne0$ and every digit
denominator $L_i$ is a nonzero power of $q$; its exponent is one of
$t,t+1,y+1,x$, all at most $x$.
For $L_i v_i=M$ and $v_i U_i=M Q_i+R_i$, the nonzero integer $v_i$
divides $R_i$.  Hence $r_i=R_i/v_i$ is an integer and
$U_i=L_iQ_i+r_i$.  The unchanged common norm
\[
 M^2-1-4\sum_i R_i^2=a^2+b^2+c(c+1)
\]
gives $|r_i|<|L_i|/2$ without assuming signs for $L_i,v_i,M$.
Any congruent canonical tail of magnitude $<|L_i|/32$ must therefore
coincide with $r_i$: their difference is a multiple of $L_i$ of
magnitude less than $(1/2+1/32)|L_i|<|L_i|$.

\subsection{Authenticating the binomial at a signed radix}
Put $n=R+t$ and $C_j=\binom nj$.  The numerator request is exactly
$U=(q+1)^n$.  For $0\le j\le t$,
\[
 C_j\le(R+1)^t\le(F^2+1)^t=:V.
\]
The standard uniform binomial bound follows by Pascal induction and the
last-edge product bound.  Here $R=F\ge1$; the mod-four-only controls
instead retain $R=F^2$ and satisfy the same upper bound.
The symbol $V$ is proof notation, not a new power slot.
For $F\ge7$, $F^2+1\le2^{F-1}$ (the base case and induction are immediate).
Equation~\eqref{eq:r20-q-size} gives
\[
 V/Q\le2^{-F-t-2},\qquad V/(Q-1)<1/32.
\]
For either sign of $q$, the signed low polynomial
\[
 r^0=\sum_{j=0}^{t-1}C_jq^j
\]
has $|r^0|<Q^t/32$ and is congruent to $U$ modulo $q^t$.
The short-lift argument identifies the digit remainder with $r^0$.
Dividing the digit equation by $q^t$ and reducing modulo $q$ now gives
$B\equiv C_t\pmod q$.
But $0<C_t<Q/32$ and $|B|\le Q/4$, so $|B-C_t|<Q$:
\[
 B=\binom{R+t}{t}>0.
\]
Since $H>0$, the self-radix equality then proves $q=(BH)^x>0$.
The initially negative branch has been excluded by arithmetic.
The request $N=B^x$ now has its ordinary positive nonunit base meaning.
Also $B\ge R+t>R+1$, so the raw index equation excludes $C_0=0$.

The remaining range and local-congruence arguments are unchanged.
Their canonical tails are $<L_i/32$; all fixed front indices and local
majorants are no larger after replacing $F^2$ by $F$.
The strict endpoint budget $x\ge2+(t+1)F$ is preserved, including the
derived-label bound.  Applying the two-sided binomial selector theorem
using~\eqref{eq:r20-linear-radius} recovers every local C222 row.

\subsection{Recovering the proof and constructing all witnesses}
The recovered beta rows give accumulators $S_0,\ldots,S_t$ and updates
\[
 S_{i+1}(1+c_iF)=S_i(1+p_iF)(1+q_iF),
\]
whose weights are positive.  If $y=F$, the first endpoint would give
$S_0=0$, while the last gives $S_t=1$; every update propagates zero,
a contradiction.  Thus $y>F$, the initial remainder is genuinely
$F+1$, the final one is $1$, and all accumulators are positive.
Multiplication of the updates yields the inherited label-product identity.
The two sides have equal nonunit factor counts $d\le t$: if they differed,
the side with more factors would be at least $(F+1)^{d+1}$, while the
other is at most $t^d(F+1)^d$, contradicting $F+1>t^t$.
Each side is consequently the value at $F$ of a product of at most $t$
factors $1+jX$, $1\le j\le t$.  Its coefficient sum is at most
$(t+1)^t<F$.  Reduction modulo $F$ at the least nonzero degree shows
that a nonzero difference cannot vanish at $F$.
The products are identical polynomials; their distinct roots $-1/j$
recover every label multiplicity.  The retained Horn-closure argument
therefore gives a finite C222 refutation and hence a finite ZFC refutation.
This proves soundness.

Conversely, begin with a finite C222 refutation and choose a sufficiently
large padded horizon $t=16+h^2$.  Set $F=[2(t+1)]^t$, $R=F$, and
$B=\binom{F+t}{t}$.  The two-sided selector theorem applies by
\eqref{eq:r20-linear-radius}.
Let $O$ be the odd part of $t!$.  Choose an arbitrarily large odd multiple
$a_\beta$ of $O$, and put $y=F a_\beta^2$.
Since $v_2(F)\ge t>v_2(t!)$, this is an even multiple of $t!$ and can
exceed $F$ and every accumulator.  The beta CRT works as before; its
moduli are odd, so choose a sufficiently large even beta code $x$.

The smaller-radius binomial is still odd:
$v_2(F+i)=v_2(i)$ for $1\le i\le t$, since
$v_2(i)<t\le v_2(F)$.  Thus $v_2(B)=0$ and $N^2=B^{2x}$ is odd.
The all-primes CRT construction provides the actual representatives.
Adding $N^2$ permits their parities to be chosen while keeping the old
bounds.  Choose the dominant $a$ odd and every other representative even;
then $a,b$ have the required opposite parities, the nonzero semantic
anchor is even, and $H=2+\sum a_\nu$ is odd.
The changed front index satisfies $R+t\le B<B^2<N^2$, so every retained
principal-index and frame bound remains valid.
Now $q=(BH)^x$ is an odd square, hence $q\equiv1\pmod4$, and
$z=(q-1)/4$ supplies the new affine radix constraint.
Canonical normalized tails are $R_i=(M/L_i)r_i^0$ with
$0\le R_i<M/32$; the strict common norm gap has the usual ternary
representation.  The unchanged Pell completion supplies all other
witnesses.  Factorials are used only to select a divisible beta parameter,
not as an authenticated unknown.  This proves completeness.

\subsection{Literal source edits, exact inverse certificates, and limits}
In both pinned source equations $F=x_3$, $R=x_4$, and row zero is
$R-F^2$.  The source-interface audit checks that all previous factoring,
projection, translation, and row-basis steps retain this row and these
coordinates without using the row as a premise in an elimination.
Replacing it by $R-F$ therefore implements exactly the semantic radius
change.  The centered stored radix has actual value
$q=\text{stored}_q-1$; its old row is
$\text{stored}_q-2-4z-4z^2$.  The new row is
$\text{stored}_q-2-4z$, so its affine restoration is
$\text{stored}_q=4z+2$, not $4z+1$.
The second restoration is $R=F$.

These semantic edits preserve \emph{zero existence}, by the two ordinary
arithmetic directions above.  They are not claimed to be bijections
between all predecessor witness tuples.  Only after the source edits do
the two monic affine eliminations give exact restriction/restoration
bijections.  Every later row operation $R_i\leftarrow R_i+kR_j$ has the
integer inverse with $-k$ and preserves the common zero set, though it
generally changes the sum-of-squares polynomial.

The independent checkers import no new producer, verify all restored
coordinates and pulled-back residuals, reverse every frozen row operation,
collect every quartic coefficient, and check degree, coordinate liveness,
and all reported metrics.  Deliberate corruptions must fail, under both
normal Python and Python~\texttt{-O}.  These are finite artifact checks,
not replacements for the unbounded arithmetic proof.
The complete immediate predecessor, its ordinary proofs and calibration
interface sources, is retained under \texttt{inherited\_factorial\_free}.
The named logical interfaces are
\begin{quote}\small\ttfamily
KEndpoint.c222\_iff\_zf\_refutation\\
ConstructibleChoice.zf\_refutation\_iff\_zfc\_refutation.
\end{quote}
Their composition yields the stated equivalence with $\operatorname{Con}(\mathrm{ZFC})$.
No global minimum, proof of consistency, new proof-assistant closure, or
literature-wide priority is claimed.

\section{Inverse-Lucas correction: 85 operations in ten positive variables}
\label{sec:inverse-lucas}

\paragraph{Result and scope.}
The literal polynomial \texttt{positive/equation\_bV.json} has ten
independently strictly positive integer unknowns, exact degree \(304\),
and \(85\) live binary arithmetic operations.  The newest inspected
Lucas-correction parent has the same degree and \(87\) operations.
The square-weight control has \(84\) operations and degree \(410\);
it is dominated by known earlier points and is not promoted as a frontier
improvement.  The complete indexed BFS943/radix67 tuple is unchanged,
with largest fixed port of \(592{,}884\) bits.

Composed with the explicitly retained ordinary native/refutation
calibration, the proof below gives
\[
 \mathrm{ZFC}\vdash
 \bigl(\neg\exists\boldsymbol{x}\in\mathbb Z_{>0}^{10}\;
 P(\boldsymbol{x})=0\bigr)\longleftrightarrow\operatorname{Con}(\mathrm{ZFC}).
\]
This neither proves consistency nor asserts a zero.  It is an ordinary
mathematical proof with exact finite artifact checks, not a new end-to-end
proof-assistant object or a global-minimum claim.

The new mechanism replaces the forward rank correction by an inverse
Lucas correction at the existing abscissa \(k\).  This saves two
multiplications and removes the exceptional negative-rank transport cases.
A genuine obstruction at the prime \(3\) means that the old denominator
proof cannot simply be reused.  Projection must precede full denominator
recovery; a new absolute bound handles a possibly negative corrected quotient
in the positive lift.  The second-order expansion and the projection-first
strategy themselves are explicitly inherited from the inspected Lucas
Corrections and Cofinal Witnesses archives at commit
\texttt{7a8da414376a80b90ecf799f37ceb26fb9ac674e}.

\subsection{Complete polynomial}
The only unknowns are
\[
 J,F,\alpha,i,T,w,s,n,\eta,\zeta\in\mathbb Z_{>0}.
\]
The fixed integers \(\beta,K_0,M_C,M_F\) are the whole supplied tuple,
with \(\beta\ge31\) odd, \(K_0>0\), and
\(0<M_C<\beta<M_F<2\beta\).  Define polynomial abbreviations:
\begin{gather*}
 h=\beta J,\quad q=h+1,\quad C=q-F-\alpha,\\
 X=qw,\quad Y=qs,\quad k=\eta+\zeta,\quad c=kY+\eta,\\
 A=Y(X+1)+1,\quad \Delta=A^2-1,\quad H=4A-5,\\
 v=XY^2,\quad \ell=vk=(XY)(kY),\quad E=1+\ell(\ell+k),\\
 R=[h(q-F)+\alpha](q^2-1)+(M_C+qM_F)J,\\
 B=(K_0+w)C-F,\quad U=ic^2,\quad a=\Delta U^2E,\\
 F_0=1+a,\quad G=EF_0,\quad F_s=1+4aF_0,\quad Q=\Delta(F_s-1),\\
 b=cT-R,\quad V=F_sb-c,\quad K=X(3c-2X)+2,\\
 e_6=6(k-R-1)-\ell(k^2-4),\quad D_6=6q^6,\quad M=hD_6H.
\end{gather*}
Use \(L=bV\) for the principal equation, or \(L=V^2\) for the control,
and put
\begin{align}
 Z&=hD_6\bigl[H(n-LB)+K\bigr]
       -H\bigl[D_6B+h(e_6-6q^4s)\bigr],\label{eq:inverse-Z}\\
 P&=Q\bigl(Z^2-M^2GV^2\bigr)+(M^2G-Z^2).\label{eq:inverse-P}
\end{align}
The single equation is \(P=0\).  All displayed powers have fixed
nonnegative integer exponents.  The DAG does not compute \(v\)
separately: it reuses \(\ell=(XY)(kY)\).  There are no division gates,
variable exponents, sequence functions, or additional quantified variables.

For proof only define
\begin{gather*}
 e_*=\frac{e_6}{6}=k-R-1-\frac{vk(k^2-4)}{6},\quad
 \rho=\frac{e_*}{q^6},\quad\sigma=\frac{s}{q^2},\\
 t=\frac{B}{h},\quad W=1+hL,\quad
 \kappa=\frac{K}{H},\quad z=\frac{Z}{M}.
\end{gather*}
Without assuming any integrality of these quotients,
\begin{equation}
 z=n+\kappa+\sigma-Wt-\rho.\label{eq:inverse-carrier}
\end{equation}

\subsection{Root extraction and uniform signed transport}
The unchanged root and folding arguments are retained in
\texttt{inherited\_lucas/PROOF.md}, Sections 2--3, with their ordinary
Pell sources.  We state their premises and conclusions to specify the
interface precisely.  Positive inputs give \(q\ge32\), \(X,Y\ge q\),
\(k\ge2\), \(c>k\), \(A>4\), \(U\ge c^2\), \(F_s>2c\),
and \(Q>c^2\).  Moreover
\[
 F_s=(1+2a)^2,\qquad Q=4\Delta^2U^2G,\qquad\gcd(E,1+a)=1.
\]
At a zero, prime valuations in
\((Q-1)z^2=G(QV^2-1)\) give \(G=g^2\), \(E=\tau^2\),
\(F_0=f_0^2\), with \(g=\tau f_0>0\) integral.  Rational adjacent-Pell
descent makes \(y=|z|/g\) integral, and
\[
 QV^2-(Q-1)y^2=1,\quad z=\varepsilon gy,\quad
 \varepsilon\in\{1,-1\}.
\]
Writing \(f=1+2a\), \(I=2Ug\), we have
\(f^2-\Delta I^2=1\), \(Q=(\Delta I)^2\), and \(c^2\mid I\).
Signed all-index folding gives
\[
 c=\psi_A(p)\quad(p>0\text{ odd}),\qquad
 k=2\psi_{1+2v}(m),\quad m\ge2,
\]
where \(\psi_a(0)=0\), \(\psi_a(1)=1\), and
\(\psi_a(j+2)=2a\psi_a(j+1)-\psi_a(j)\).
In particular \(k\ge8v+4\), \(c>8XY^3\ge8q^4\), \(c>X\), and
\(K>0\).

The inherited exclusions of main and auxiliary index three use only
\(q\mid X,Y\), not the rank correction.  They give, for either weight,
\begin{equation}
 b\ne0,\quad \operatorname{sign}V=\operatorname{sign}b,
 \quad |V|>|b|,\quad L>0,\quad gy<L/2,
 \quad {|b|\over q^6}<L/2.\label{eq:inverse-abs}
\end{equation}
For \(L=bV\), the main index-three remainder would force
\(2Y-X\in\{1,2\}\), impossible modulo \(q\).  The two signed auxiliary
index-three possibilities would give
\[
 F_s(b-4\Delta)=c-(4\Delta+3),\qquad
 F_s(b+4\Delta)=c+(4\Delta+3),
\]
whose right sides lie strictly between \(0\) and \(F_s\).
Thus the auxiliary index is at least five, giving \(|b|>4g\) and
\(gy<3L/8\).  Neither sign branch is discarded.

The new correction has an advantage before any sign of \(R\) is known:
\(k\ge2\) makes \(vk(k^2-4)/6\ge0\).  Therefore
\[
 e_*\le k-R-1<b,\qquad \rho<|b|/q^6.
\]
If \(B\le-1\), then \(Wt<-L\), and
\eqref{eq:inverse-carrier}--\eqref{eq:inverse-abs} force
\[
 n=\varepsilon gy-\kappa-\sigma+Wt+\rho
   <L/2-L+|b|/q^6<0,
\]
a contradiction.  This handles all negative \(R\) and \(R=0\)
uniformly.  Hence \(B\ge0\), so \(C>0\).  Only now is the inherited
mask bound \(3q<R<q^4-2\) invoked.  Then \(c>2R\) and \(b,V>0\);
the normalized-rank and signed auxiliary-index theorems give \(p=R\).
Thus \(R\) is positive odd; mask parity gives even \(q\) and odd \(J\).

\subsection{Inverse-Lucas authentication at all primes}
The inherited, universally valid finite Lucas expansion gives
\begin{equation}
 k=2m+8vC_m+32v^2Z_m,
 \quad C_m=\binom{m+1}{3},\quad Z_m\in\mathbb Z_{\ge0}.
 \label{eq:inverse-expansion}
\end{equation}
Binomial notation occurs only in the proof.  Since \(k=2u\) is even,
\(k(k^2-4)/6=4u(u-1)(u+1)/3\) is integral.  Thus \(e_*\) is
integral at a zero.  Clearing the integral carrier gives
\(6q^6\mid hH(e_6-6q^4s)\); cancelling the ordinary factor six and
using \(h=q-1\), \(H\equiv-1\pmod q\), gives
\begin{equation}
 q^6\mid e_*-q^4s,\qquad q^4\mid e_*.
 \label{eq:inverse-q4}
\end{equation}
No modular inverse of six is used.

Put \(t_0=2m\), \(d=2m-R-1\), \(A_m=C_m+4vZ_m\), and
\(S_m=k^2+kt_0+t_0^2-4\).  The exact inverse identity is
\begin{equation}
 3(e_*-d)=4v^2(24Z_m-A_mS_m).
 \label{eq:inverse-identity}
\end{equation}
Indeed \(t_0(t_0^2-4)=48C_m\), \(k-t_0=8vA_m\), and
\(k(k^2-4)-t_0(t_0^2-4)=(k-t_0)S_m\).
Since \(v=q^3ws^2\), the right side is divisible by \(q^6\).
For each prime \(\lambda\mid q\), with \(a_\lambda\ge1\) its
valuation in \(q\),
\[
 \operatorname{ord}_\lambda(e_*-d)
 \ge6a_\lambda-\operatorname{ord}_\lambda(3)\ge4a_\lambda.
\]
Hence \(q^4\mid d\).  The first base \(1+2XY^2\) exceeds \(A\).
If \(m\ge R\), monotonicity would give
\(k/2\ge\psi_A(R)=c\), contrary to \(c>k\).  Thus \(1\le m<R\),
and \(|d|<R<q^4\), proving
\begin{equation}
 d=0,\qquad m=(R+1)/2.\label{eq:inverse-index}
\end{equation}

It would be false to infer full \(q^6\)-divisibility here.
For \(q=3,v=27,m=2,R=3\), the inverse error is
\(-47{,}911{,}824\equiv243\pmod{729}\).
This modular diagnostic is not a purported zero of the full polynomial.
Evenness of \(q\) alone does not remove the obstruction at three.

\subsection{Projection first, denominator second}
Because \(z=Z/M\) is integral, reducing the literal numerator
\eqref{eq:inverse-Z} modulo \(H\) gives \(H\mid6hq^6K\).
As \(\gcd(q,H)=1\),
\begin{equation}
 H\mid6hK.\label{eq:inverse-project}
\end{equation}
This uses neither integral \(\rho\) nor integral \(\sigma\), and does
not cancel three.  Write \(R=2r+1\), so \(r\ge48\).
The retained uniform ratio theorem, licensed by the authenticated
indices, gives
\[
 Y>X^r/3,\qquad X^{r+1}<3Y(X+1)<H.
\]
Modulo odd \(H\), the scaled recurrence gives
\(3\cdot2^{R-1}\psi_A(R)\equiv4^R-1\), because \(4A\equiv5\).
Consequently \(2^{R-1}K\equiv-(X-2^R)(2^RX+1)\pmod H\).
Equation \eqref{eq:inverse-project} therefore implies divisibility
of this whole product by \(H\), with multiplier \(6h\).

For an explicit height check let \(u_0=2^R=2\cdot4^r\).
Since \(6<q^3\) and \(q^3h<X^4\),
\begin{align*}
 6h\,|(X-u_0)(u_0X+1)|
 &<X^4(X+u_0)(u_0X+1)\\
 &\le4\cdot4^rX^6+8\cdot16^rX^5<X^{r+1}<H.
\end{align*}
After division by \(X^{r+1}\), the two terms decrease with
\(X\ge32,r\ge48\); their sum at \((32,48)\) is
\(2^{-117}+2^{-25}<1\).  Hence the whole divisible product is zero,
and its other factors are positive.  Thus \(X=2^R\); the same universal
projection then gives integral positive \(\kappa=K/H\).

Now \(q\mid X\) implies \(q=2^a\), \(a\ge5\).
Equation \eqref{eq:inverse-identity}, with \(d=0\), finally gives
\(q^6\mid e_*\), because three is a unit modulo \(q\).
Thus \(\rho\) is integral.  The carrier says
\(\sigma-WB/h\) is integral.  Multiplying by \(q^2h\) and reducing
modulo \(q^2\) gives \(q^2\mid s\).  Hence \(q^3\mid Y\),
\(\sigma\) is integral, and \(h\mid WB\).  Since \(W\equiv1\pmod h\),
\(h\mid B\), so \(t\) is a nonnegative integer.

The fixed-minus parity, exact half-binomial rounding, and fresh original
native reconstruction now have their full retained premises.
That source uses base \(Y(X+1)+2\), not the auxiliary base \(A\).
The unchanged whole-tuple calibration produces a finite ZFC refutation.

\subsection{Constructive positive completeness}
A finite refutation supplies the retained canonical outer data
\(q=2^a\), \(X=2^R=qw\), \(R>3q\), \(R\equiv3\pmod4\),
\(q^3\mid Y\), \(C>0\), and \(t=B/h\ge0\) integral.
Set \(s=Y/q>0\), and form
\[
 c=\psi_A(R),\quad k=2\psi_{1+2XY^2}((R+1)/2),\quad
 \tau=\chi_{1+2XY^2}((R+1)/2),\quad
 \eta=c-kY,\quad\zeta=k-\eta.
\]
The exact ratio theorem gives \(0<\eta<k\), hence both last inputs
are positive.  Equation \eqref{eq:inverse-identity} and dyadic \(q\)
give integral \(\rho\).  Projection gives positive integral \(\kappa\),
and \(\sigma=Y/q^3\) is integral.

A positive period of the Pell matrix modulo \(c^2\tau\) gives a seed
\((f_0,I_0)\) with \(c^2\tau\mid I_0\).
Set \(i=I_0/(c^2\tau)>0\).  Then \(a=\Delta I_0^2=f_0^2-1\),
\(F_0=f_0^2\), and \(G=(\tau f_0)^2\).
The doubled pair \(f=2f_0^2-1\), \(I=2f_0I_0\) has exactly the
required definitions.  The retained auxiliary-completeness theorem gives
positive \(V,y\) and the positive integer
\(T=(V+c+Rf^2)/(cf^2)\), the actual unshifted input.

The new \(\rho\) may be negative; the old nonnegative-quotient argument
is not reused.  Since \(k\ge R+1\) and \(vk^2>6\),
\[
 |\rho|\le {k-R-1+vk(k^2-4)/6\over q^6}
       <{k+vk^3/6\over q^6}<vk^3/3.
\]
For every permitted \(i\ge1\),
\[
 g>\sqrt\Delta\,c^2E>c^2v^2k^2>vk^3,
\]
so \(|\rho|<g/3\).  Also \(\sqrt\Delta>A-1=Y(X+1)>X\) and
\(cE>4\), hence \(g>4Xc\).  As \(\kappa<Xc\) and
\(\sigma\le Y<Xc\), we have \(\kappa+\sigma<g/2\).
For either positive weight put
\[
 n=gy+(1+hL)t+\rho-\kappa-\sigma.
\]
This is integral, and \(y\ge1\), \(t\ge0\) imply
\(n>g-g/3-g/2>0\).  Its carrier is \(z=gy\), so \(P=0\).
All ten independent coordinates have been constructed strictly positively;
no inverse positive shear or unproved sign is used.  Every permitted
period seed works, without additional cofinal enlargement.

\subsection{Exact costs and verification}
Deleting the incoming gate \(v=(XY)Y\), restoring the sharing
\(\ell=(XY)(kY)\), and replacing its four-gate rank cubic by
\(k^2\), \(k^2-4\), \(\ell(k^2-4)\) saves two multiplications.
The principal has \(51\) multiplications, \(18\) additions and
\(16\) subtractions.  All inputs, ports and rows are live.

Put \(q_0=\beta J\), \(k_0=\eta+\zeta\),
\(C_0=q_0-F-\alpha\), and
\(B_0=k_0sT-q_0^2(q_0-F)\).  The \emph{complete} highest homogeneous
forms are
\begin{align*}
 [P_{bV}]_{304}
 &=-9\cdot2^{12}q_0^{128}w^{34}s^{78}i^{14}k_0^{44}B_0^2,\\
 [P_{V^2}]_{410}
 &=9\cdot2^{16}q_0^{166}w^{46}s^{104}i^{20}k_0^{60}B_0^4C_0^2.
\end{align*}
Their full expansions have respectively \(274\) and \(3750\) terms.
Nonzero coefficients at maximal \(T\)-power establish the exact degrees
for the actual \(\beta>0\), not merely a formal upper bound or a ray sample.

The producer-independent checker covers all actual rows with \(37\)
exact free-polynomial cuts per graph, verifies full leading forms and
every coefficient byte, reconstructs the four fixed numerals, and rejects
both alternative operations at every gate plus a missing-correction
mutation: \(171\) rejections for the principal and \(169\) for the control.
Separate finite recurrence controls include the obstruction at three;
they do not replace the unbounded arithmetic proof.
Normal and optimized Python receipts agree byte for byte.
\begin{verbatim}
python3 -B positive/build.py --check
python3 -B positive/verify.py
python3 -O -B positive/verify.py
\end{verbatim}

\subsection{Whole complemented-label radix47 sidecar}
The same two free-port row lists are separately specialized to the
\emph{entire} frozen complemented-label radix47 tuple.  The principal
remains \(85/304\), but its largest fixed port is \(415{,}904\) bits and
the four raw numerals occupy \(207{,}214\) bytes rather than \(295{,}388\).
This reuses the incoming encoding and controller; it is not a new radix
theorem and is not the repository's conditional three-anchor radix47 layout.
The canonical radix67 equations remain separate.

The new arithmetic is uniform in the same scalar, source-shift mask and
population interface already checked for the whole47 tuple.  In soundness
it recovers \(C>0\), \(X=qw=2^R\), \(q^3\mid Y\), the native
half-binomial, and \((K_0+w)C-F=t(q-1)\), \(t\ge0\).
These are exactly the retained whole47 bridge's premises, with
\(W_{\rm native}=w\) and source transport variable \(t+1\).
Its fresh Pell witnesses are constructed at \(Y(X+1)+2\), before the
actual complemented-label decoder is applied.  Conversely, the actual
whole47 controller supplies dyadic \(q\) and the same outer interface;
choosing \(w=W_{\rm native}\), \(s=Y/q\), and applying the explicit
positive construction above gives a new zero.  No old67 decoder or
controller is silently substituted.

The complete ordinary native bridge, incoming P66 equation, all four
payloads, \(3297\) coefficients, \(51\) native target identities, and its
\(37\) pinned dependencies are retained under
\path{positive/radix47/inherited/}.  The sidecar checker requires identical
free-port rows and an exact whole-tuple binding, reconstructs every numeral,
and repeats the cuts, leading-form and mutation checks.  Its normal and
optimized receipts agree.  The default main checker includes this sidecar;
its literal binding can also be reproduced independently:
\begin{verbatim}
python3 -B positive/radix47/build.py --check
python3 -B positive/radix47/verify.py
\end{verbatim}

\section{Thirteen integers at degree 152 in 119 operations}
\label{sec:root119}
The literal polynomial \texttt{integer13\_degree152\_ops119.json} has thirteen
unrestricted integer inputs, exact degree $152$, and $119$ live binary operations:
$71$ multiplications, $24$ additions and $24$ subtractions. Squares and scalar
multiplications are charged. This strictly improves the incoming $13/152/126$
evaluator. Against the earlier $13/156/117$, it trades two operations for four
degrees. The new whole polynomial is exactly the incoming one under
$n_{\rm old}=n_{\rm new}-k$; this is a polynomial automorphism, not a numerical
or generic-parameter comparison. All six supplied coefficient integers are
unchanged. Intermediate registers below are abbreviations, not unknowns.

The ordinary signed-index and native/compiler/refutation theorem is inherited
from \path{roots/inherited/PROOF.md}, Sections 1--5, with its full retained
source packet and finite numeral recipe. Thus ordinary ZFC proves that universal
integer nonvanishing is equivalent to syntactic $\operatorname{Con}(\mathrm{ZFC})$.
This does not assert consistency, global minimality, a new native encoding or
a new end-to-end proof-assistant certificate. The incoming improvement from
$XY$ to $XY^2$ in the first-index modulus and from $q^8$ to $q^6$ in the budget
is credited to that source, not to the present evaluator change.

\subsection{The full division-free polynomial}
The thirteen inputs, all in $\mathbb Z$, are
\[
 j,C,z,h,r,v,w,T,c,n,y,p_1,p_2.
\]
The complete fixed ports are $\beta,K_0,M_C,M_F,A_c,B_c$, with $M_F$ already
shifted. In particular $\beta=2^{1064921}-1$, $A_c=2^{1113}$ and
$B_c=2^{666104}$. Their full bytes are in the component's
\texttt{inherited/coefficients}, not unspecified parameters. Define
\begin{gather*}
 J=6j+1,\quad q_m=\beta J,\quad q=q_m+1,\\
 W=12w+4,\quad X=qW,\quad Y=v,\\
 F=(K_0+W)C-zq_m,
 R=[q(q-F)-C](q^2-1)+(M_C+qM_F)J,\\
 E=XY,\quad a=E+Y,\quad H=4a+3,\quad \Delta=a^2+H,\\
 \ell=EY,k=R+1+h\ell,\eta=c-Yk,L=\ell k,\\
 I=(2r+1)c^2,D_s=\Delta I^2,F_s=1+D_s,Q=\Delta D_s,\\
 V=F_s(cT-R)-c,\quad q_3=q^3,M=2q_3,U=MH,\\
 K_{\rm lo}=X(2X-3c)-2,
 D_{\rm err}=(A_cF-q)^2+(B_cC-q)^2,\\
 B=(Yk)^2(q^2-D_{\rm err})-q^2[4(Y\eta)^2+(kq^2)^2],\\
 G=4(B-p_1^2-p_2^2)-3,D_t=1+L(L+k),\\
 g_a=U^2G,d_b=4U^2D_t,f_c=16U^2F_s,\\
 Z_c=U(n+2L)+MK_{\rm lo}+H(Y-q_3),\\
 g_c=g_a+f_c,Z_2=Z_c^2,g_{af}=g_af_c,\\
 K=(Z_2-g_c-d_b)^2-4(g_{af}+d_bg_c),\\
 N=K^2-64g_{af}d_bZ_2,u_2=(V+y)^2,\\
 E_a=u_2-Q(u_2-V^2),P=(N+1)E_a-1.
\end{gather*}
The auxiliary identity
\begin{equation}\label{eq:roots-aux119}
 E_a=V^2-(Q-1)y(2V+y)
\end{equation}
is exact. Reusing $u_2$ saves one binary operation.

\subsection{Whole-polynomial identity and inverse chart}
Write, for proof only,
\[
 Z=Un+MK_{\rm lo}+H(Y-q_3),\quad b=2L+k,\quad
 A(z)=(z+Uk)^2+4U(Lz-U).
\]
The incoming norm uses
\begin{align*}
 K_{\rm old}&=(A(Z_{\rm old})-g_a-f_c)^2
                 -4[g_af_c+d_b(g_a+f_c)],\\
 N_{\rm old}&=K_{\rm old}^2
                 -64g_ad_bf_c[A(Z_{\rm old})+d_b].
\end{align*}
Under $n_{\rm old}=n-k$ its numerator becomes $Z_{\rm old}=Z-Uk$, and
\begin{align}
 A(Z-Uk)&=Z^2+4ULZ-4U^2(1+Lk),\label{eq:roots-cut119}\\
 A(Z-Uk)+d_b&=(Z+2UL)^2=Z_c^2.\label{eq:roots-square119}
\end{align}
Therefore $K=K_{\rm old}$ and $N=N_{\rm old}$ identically. Together with
\eqref{eq:roots-aux119}, this proves
$P_{\rm new}(n)=P_{\rm incoming}(n-k)$ as whole polynomials. Since $k$
is independent of $n$, the inverse is $n_{\rm new}=n_{\rm old}+k$ on every
integer assignment. Only the degree-five shift $k$ is removed; removing the
degree-nine $2L$ as well would destroy the degree-$152$ cancellation.

\subsection{All signed roots and all boundary cases}
We give direct maps to the fifteen-input parent as an additional unbounded
check. Unconditionally, $q$ is even and nonzero, because $\beta>1$ is odd and
$\beta(6j+1)=-1$ is impossible. Thus $|q|\ge2$, $|W|\ge4$, $|X|\ge8$,
and $X$ is even. Since $a=(X+1)Y$ is zero or has absolute value at least $7$,
\[
 \Delta=(a+2)^2-1\ge3.
\]
$H$ is odd and nonzero, so $U\ne0$. Also $\ell=XY^2$ is even,
$\ell(\ell+1)\ge0$, and
\[
 D_t=1+\ell(\ell+1)k^2\ge1\text{ is odd},\qquad F_s\ge1.
\]
$G\equiv1\pmod4$ is odd and nonzero. These facts hold before any guard or
zero assumption. Since $N$ is a square modulo $4$, $N\ne-2$. At $P=0$ the
integer factors $N+1,E_a$ are units; the first cannot be $-1$. Thus
$N=0,E_a=1$, and the converse is immediate.

Put $\rho=Z_c/U\in\mathbb Q$. The universal complex identity is
\begin{equation}\label{eq:roots-factors119}
 \frac{N}{U^8}=
 \prod_{\varepsilon\in\{\pm1\}^3}
 (\rho-\varepsilon_1\sqrt G-2\varepsilon_2\sqrt{D_t}
                           -4\varepsilon_3\sqrt{F_s}).
\end{equation}
It follows by multiplying all eight factors. If $G<0$, its square root
is the sole nonreal term of each factor, with nonzero imaginary part; therefore
$N=0$ is impossible. Hence $G>0$. The case $G=0$ is excluded by oddness;
$D_t,F_s$ are strictly positive, including their possible value $1$.

Write the three radicands as $u_i^2d_i$ with $u_i>0$ integral and $d_i$
squarefree. Distinct squarefree radicals, including $\sqrt1$, are independent
over $\mathbb Q$: in the multiquadratic field generated by all their primes,
character averaging of the sign-change automorphisms projects onto each radical.
In any vanishing factor of \eqref{eq:roots-factors119}, group terms by squarefree
class. An irrational group containing $G$ has an odd coefficient plus even
coefficients. A group without $G$ but containing $D_t$ has a coefficient of
exact $2$-adic valuation one plus at most a multiple of four. A group containing
only $F_s$ is nonzero. None can cancel; $u_i$ for $G,D_t$ are odd. An even
squarefree class can contain only $F_s$ and is excluded likewise. This explicitly
handles shared classes; no generic distinctness assumption is used.

Every class is therefore $1$, and the signs of the zero factor give integers
\begin{equation}\label{eq:roots-integrals119}
 g^2=G,\quad\tau^2=D_t,\quad f^2=F_s,\qquad
 \rho=g+2\tau+4f.
\end{equation}
In particular $\rho$ is odd and nonzero. No division by $\rho$ was used.
All eight sign choices are allowed; $f$ and $\tau$ were not assumed positive.

The fifteen-input kernel retains the common coordinates but has additional
$f,p_3$, and residuals
\begin{align*}
 A&=(Z+Uk)^2+4U(LZ-U),& B_s&=f^2-\Delta I^2-1,\\
 C_a&=V^2-(Q-1)y(2V+y)-1,&&
 G_b=B-1-p_1^2-p_2^2-p_3(p_3+1),
\end{align*}
where its $Q=\Delta(f^2-1)$ and $V=f^2(cT-R)-c$.
At a new zero restore
\[
 f_{\rm old}=f,\quad p_{3,\rm old}=(g-1)/2,\quad
 n_{\rm old}=n-k-g-4f.
\]
All are integral. Its centered carrier satisfies
\[
 Z_{\rm old}/U+(2L+k)=\rho-g-4f=2\tau.
\]
The identity $A=(Z+U(2L+k))^2-4U^2D_t$ gives $A=0$.
The equation $f^2=F_s$ gives $B_s=0$ and identifies both versions of $Q,V$;
then $C_a=0$. Finally
$p_{3,\rm old}(p_{3,\rm old}+1)=(G-1)/4=B-p_1^2-p_2^2-1$, so $G_b=0$.
This is an actual parent zero.

Conversely, from any parent zero,
$(Z_{\rm old}/U+2L+k)^2=4D_t$. A rational square root of an integer is
integral by reducing its fraction to lowest terms, and this root is even
because its square is divisible by four. Write it $2\tau$. Set
$g=2p_{3,\rm old}+1$ and
\[
 n_{\rm new}=n_{\rm old}+k+g+4f_{\rm old}.
\]
Delete $f_{\rm old},p_{3,\rm old}$ and retain all other inputs. The old guard
and strong row give $g^2=G$ and $f_{\rm old}^2=F_s$. The new carrier is
$\rho=g+2\tau+4f_{\rm old}$, so a factor of \eqref{eq:roots-factors119}
vanishes and $E_a=1$. Hence $P=0$. These constructions apply to every parent
zero and add no new simultaneous Pell or canonical-witness existence claim.

\subsection{Exact degree and complete leading cancellations}
Naive max/sum degree propagation on the literal DAG gives only an upper bound
$172$. Both complete cancellation cuts below are essential.
With all ports of degree zero, put
\begin{align*}
 q_0&=6\beta j,&X_0&=12q_0w,&Y_0&=v,\\
 F_0&=12wC-q_0z,&R_0&=-q_0^3F_0,&\ell_0&=X_0Y_0^2,\\
 k_0&=R_0+h\ell_0,&a_0&=X_0Y_0,&U_0&=8q_0^3a_0,\\
 T_0&=A_c^2Y_0^2F_0^2+4q_0^2Y_0^4+q_0^6.
\end{align*}
The entire budget top is $[B]_{16}=-k_0^2T_0$; in particular the
$h\ell_0$ part of $k_0$ must not be dropped. The scaled radicands have degrees
$28,30,24$. Although $Z_c^2$ has degree $30$, the complete identity
\eqref{eq:roots-cut119}--\eqref{eq:roots-square119} shows that $Z_c^2-d_b$
has degree $26$, with top $-4U_0^2\ell_0k_0^2$. Consequently the first square
in $K$ has degree $56$, while $-4d_bg_a$ has degree $58$, and
\begin{align*}
 [K]_{58}&=64U_0^4\ell_0^2k_0^4T_0,\\
 [N]_{116}&=4096U_0^8\ell_0^4k_0^8T_0^2.
\end{align*}
The outer subtracted product has degree $112$ and cannot cancel this square.
Independently, with formal weights
$(Z,U,L,k,G,F_s)=(7,6,9,5,16,12)$, the checker expands the entire norm
and obtains the full top $256U^8L^4G^2$ of weight $116$.

The second deliberate cancellation is $u_2-V^2=2Vy+y^2$ of degree $18$,
although $u_2$ itself has degree $34$. It gives
\[
 [E_a]_{36}=32a_0^6r^4c^8R_0y.
\]
The product of these displayed complete tops is the entire degree-$152$
output top. In particular, on the authentic coefficient slice,
\[
 [j^{90}w^{18}v^{22}z^9r^4c^8y]P
      =2^{167}3^{108}\beta^{90}>0.
\]
The exponents sum to $152$. The symbolic extraction first sets $C=h=0$,
discarding only terms with variables absent from the target, and then selects
the $q_0^6$ term in both copies of $T_0$. Thus the degree is attained,
not merely a syntactic upper bound or a generic-port assertion.

\subsection{Finite evidence and portable closure}
The independent checker imports neither builder nor circuit implementation.
It binds all cuts of the incoming $126$-gate and new $119$-gate DAGs, checks
all topology, liveness, coverage, six fixed ports and costs, proves the whole
shear and signed-factor identities, expands both leading cancellation cuts,
and extracts the stated actual coefficient. Each of the $119$ live gates is
separately corrupted and rejected. Normal and optimized Python runs both pass.
All six raw coefficient payloads are reconstructed from the inspected retained
finite recipe in both modes. The $188$ retained files are SHA-256 pinned.
These are checks of algebra and transcription, not substitutes for the
unbounded proof above or for the inherited native/compiler theorem.

No twelve-variable breakthrough is claimed. Exploratory denominator-transfer,
primitive-carrier and strong-scale variants did not produce a verified
improvement and are not presented as theorems or new global obstructions.

\section{Reproduction and mathematical boundaries}
Every direct equation is the sum of the squares of a complete finite list
of quadratic residuals. Both the residuals and every collected coefficient
are supplied. The small-variable equations are complete straight-line
DAGs over fixed integers and precisely the listed variables. No variable
exponentiation, factorial, floor, binomial predicate or existentially
unspecified constant remains in a final polynomial.

The portable top-level replay authenticates the complete release and runs
normal and optimized Python checks in fresh disposable copies. It checks
source edits, inverse/restoration maps, entire polynomial expansions,
all circuit registers, exact leading forms, fixed numeral recipes and
actual corruption controls. Finite testing is not used to replace any
unbounded arithmetic theorem. The ordinary arguments establish soundness
and completeness in the stated domains before the finite artifacts are
used to bind those arguments to the literal equations.

The direct constructions retain the numerical C222 calculus and its
ZF/ZFC refutation calibration, without K215 simulation or Jones
universality. The compact and root branches have distinct inherited
native calibrations and coefficient tuples. A whole alternate tuple
changes the specialized polynomial; ports may never be mixed merely
because they share names or satisfy some numerical inequalities.

There are different kinds of maps here. A polynomial triangular change
of unrestricted integer coordinates is an actual lattice bijection.
The natural source's new sign-oriented factors instead require a
canonical CRT witness construction; arbitrary old natural witnesses
need not lift pointwise. The factorial-free semantic changes likewise
have separate unbounded soundness and completeness proofs. These
statements are not conflated with the universal affine restoration
certificates used later in compilation.

Reviews were performed separately within this AI-assisted workflow;
they are not outside human endorsements. No upstream repository was
modified. No global equation minimum, intrinsic arithmetic-complexity
lower bound, literature-wide priority, new Lean/Rocq endpoint or encoded
object-language ZFC derivation is asserted.
\appendix
\section{Complete fixed-coefficient recipe}\label{sec:inverse-constants}
The following finite recipe is retained from the inspected ProveIt indexed
constant specification, whose exact source and SHA-256 are recorded in
\path{inherited_lucas/parent/proof_sources/indexed_constants.tex} and
\texttt{SOURCE\_PINS.json}. It defines the entire unchanged
BFS943 tuple, not a new specialization. The symbol $K$ in the inherited
recipe is renamed $K_0$ to agree with Section~\ref{sec:inverse-lucas}; symbols
in this section are local fixed data, not additional unknowns. The supplied
standard-library reconstruction checks every resulting byte against both
literal circuits. The source shift is already included in $M_F$.
\begingroup
\providecommand{\fl}[1]{\left\lfloor#1\right\rfloor}
\providecommand{\rg}[2]{\{#1,\ldots,#2\}}
\textbf{The four fixed integers.}
Every index, set, sum, product and floor below defines a constant, not
another unknown. Put
\[
\rho=2^{67},\quad B=\rho^{8807},\quad\beta=B-1,\quad
G_m(x)=\sum_{j=0}^{m-1}x^j.
\]
Define
\[
\begin{split}
\mathcal E={}&\rg0{942}\cup\rg{983}{992}\cup\rg{1003}{1011}\\
&\cup\rg{1013}{1031}\cup\{1063,1064,1095,2195,3296\}.
\end{split}
\]
Then
\begin{align*}
M_C={}&B-1-\sum_{e\in\mathcal E\setminus\{1\}}\rho^e-2\rho^{1063},\\
M_F={}&B-1+4+(2^{51}-2)(\rho^{4393}+\rho^{5495})\\
&+\rho^{1013}G_{10}(\rho)+\rho^{1065}G_9(\rho)
  +\rho^{2195}+\rho^{3296}\\
&+(2^{65}-1)\rho^{8792}G_{13}(\rho)+(2^{19}-1)\rho^{8805},\\
K_0={}&\rho^{30}+\rho^{62}+\rho^{1100}+\rho^{2201}+\rho^{8849}\\
&+(\rho^{4393}+\rho^{5495})\Phi(1/\rho),
\end{align*}
where
\[
\begin{split}
\Phi(x)={}&(1+2^{51})G_{943}(x)
 -2^{10}\sum_{i=0}^{942}v(t_i)x^i\\
&+\sum_{(b,e,m)\in\mathcal C}2^b x^e G_m(2x),\\
\mathcal C={}&\{(10,1023,8),(18,1013,9),(27,1003,8),\\
&(35,983,9),(44,1031,1),(45,1022,1),(46,1011,1),\\
&(47,992,1),(48,1095,1),(49,2195,1),(50,3296,1)\}.
\end{split}
\]
The rational intermediate in $K_0$ cancels exactly; all four results are positive integers.

\textbf{The 943 indexed fixed tiles.}
Use the actions $u_j,v_j,d_j$ decoded below. Let
\[
\begin{gathered}
h(q,b)=2(q+1)+b,\quad
\mathcal R=\{u_j:d_j=1,\ 0\le j<430\},\\
\mathcal L=\{u_j:d_j=0,\ 0\le j<430\},\\
r_q=6+|\mathcal R\cap\rg0{q-1}|,\quad
l_q=205+|\mathcal L\cap\rg0{q-1}|.
\end{gathered}
\]
An interval ending at $-1$ is empty. The table defines $t_i$ for
$0\le i\le942$; all index ranges are fixed:
\[
\setlength{\arraycolsep}{1.4pt}\begin{array}{c|l}
\text{index }i&\text{tile }t_i=(w,n,e,s)\\\hline
0,1&(2,435,3,1),\ (5,0,0,0)\\
j+2\ (j=2q+b<430)&((1-d_j)l_{u_j},h(q,b),d_jr_{u_j},v_j)\\
432+b\ (b=0,1)&(0,h(215,b),0,h(215,b))\\
434+b\ (b=0,1)&(0,b,0,b)\\
424+2r_q+b\ (q\in\mathcal R,b=0,1)&(r_q,b,0,h(q,b))\\
424+2l_q+b\ (q\in\mathcal L,b=0,1)&(0,b,l_q,h(q,b))\\
924+j\ (0\le j<4)&(0,j+430\fl{j/2},0,438)\\
928&(2,434,2,0)\\
929&(3,436,4,4)\\
930&(4,436,4,0)\\
497+t\ (t=434,435,436)&(1,438,1,t)\\
\substack{934+3a+b-6\fl{a/4}\\a\in\{0,1,4\},\ b=0,1,2}&(a,437,b,437)
\end{array}
\]
For $t=(w,n,e,s)$ put
\[
v(t)=w+2^8n+2^{17}e+2^{25}s+127\,2^{34}.
\]

\textbf{A closed decoder for the 430 actions.}
Put $D_b(A,k)=\fl{A/b^k}\bmod b$, where
$a\bmod b=a-b\fl{a/b}$. For $0\le q<215$, $b=0,1$, $j=2q+b$:
\[
\begin{gathered}
N=\mathcal A\bmod2^{636},\quad
T_0=\fl{\mathcal A/2^{636}}\bmod2^{1675},\\
E_0=\fl{\mathcal A/2^{2311}},\\
\nu_j=D_2(N,j),\quad \sigma_j=\sum_{k=0}^{j-1}\nu_k,\\
m_q=D_2(N,430+q),\quad \tau_q=\sum_{k=0}^{q-1}m_k,\\
F_q=14(1-m_q)+m_qD_{16}(E_0,\tau_q),\\
u_j=\nu_j(1+\sigma_j)+(1-\nu_j)D_{216}(T_0,j-\sigma_j),\\
v_j=D_2(F_q,2b),\quad d_j=D_2(F_q,2b+1).
\end{gathered}
\]
Empty sums are zero. Here $(v_0,d_0,u_0)=(1,1,1)$; the constructor tile
uses this actual first target, including its state renaming.

\textbf{The integer $\mathcal A$.}
The alphabet, in order of digit values $0,\ldots,61$, is\par\noindent
{\fontsize{9}{10}\selectfont\ttfamily 0123456789ABCDEFGHIJKLMNOPQRSTUVWXYZabcdefghijklmnopqrstuvwxyz}.\par
Concatenate the following 436 characters as digits $e_0,\ldots,e_{435}$;
put $\mathcal A=\sum_{j=0}^{435}e_j62^{435-j}$.\par
\begingroup\fontsize{9}{10}\selectfont\ttfamily
BxEnhvy6NipzkWmROo0GsA3uK3twWLxscLq20UTumy5Cdu461LwD6g7RDNzF8Y\par
iYDFTpB8ehAypd7d6OsOL4APamnJeNXQypUezCmW9n4oXP6XfgmiWuAWZ8ZJ7U\par
OP7gu4kVdlGZqZHULQMar1dXNeJW4i7oNVI8usp8BlovjjZbwPCovZI328fRUq\par
4q92R81s5GjWQaQfGhkq9UYMK3rRMaBmsOIBkJqc2cj2VRIZ6mjzkBIOkW9SjZ\par
FzeA6J5u4j3sBL3SdgwXzM9mHvQivu8EJ1TBTpRhHDqZIbmYgwQ8VqTAFotl6F\par
We1YC554GKrVHsA8oUGKSA6F6k5KREBSD0Tgd7bFALY84LLbpw574BQyvdoHrl\par
0Oe35uwDf15DpT6VJWZBjSbHPpPAINYTte0eUxnTFTKBlGj8Am6sb1gLldh8WZ\par
8P\par
\endgroup

\endgroup

\section{Complete complemented-label radix47 constants}
\label{sec:inverse47constants}
This is the entire incoming complemented-label tuple, not the repository's
conditional three-anchor radix47 layout. All indices, labels and sums below
are fixed data. The finite recipe and 436-character seed are reproduced from
\path{radix47/inherited/provenance/radix47/radix47/equation.tex}; its $K$ is renamed $K_0$.
The independent reconstruction checks all 943 tiles, all 3297 coefficients,
and every byte of all four actual fixed integers. No new encoding is claimed.
\begingroup
\providecommand{\fl}[1]{\left\lfloor#1\right\rfloor}
\providecommand{\rg}[2]{\{#1,\ldots,#2\}}
\textbf{The four fixed integers.}
Every index, set, sum, product and floor below defines a constant, not
another unknown. Put
\[
\begin{gathered}
\rho=2^{47},\quad B=\rho^{8807},\quad\beta=B-1,\\
G_m(x)=\sum_{j=0}^{m-1}x^j,\quad
\Lambda=\rho^{4393}+\rho^{5495}.
\end{gathered}
\]
\begin{align*}
M_C={}&\beta+\rho-(\rho+3)\rho^{1063}\\
&-\sum_{e=0}^{3296}\fl{(D_\rho(\Phi(\rho),e)+\rho-1)/\rho}\rho^e,\\
M_F={}&2B+3+(2^{47}-2)\Lambda\\
&+\rho^{1013}G_{10}(\rho)+\rho^{1065}G_9(\rho)
  +\rho^{2195}+\rho^{3296}\\
&-\rho^{8792}+(2^{28}-3)\rho^{7694}-\rho^{7691},\\
K_0={}&\rho^{30}+\rho^{62}+\rho^{1100}+\rho^{2201}+\rho^{8849}\\
&+\Lambda\Phi(1/\rho),
\end{align*}
where
\[
\begin{split}
\Phi(x)={}&\sum_{i=0}^{942}(1+2^{10}+2^{10}w_i+2^{18}n_i\\
&+2^{27}p_i+2^{35}s_i)x^i\\
&+2^{10}x^{1003}(x^{20}+2^{17})G_8(2x)\\
&+2^{18}x^{983}(x^{30}+2^{17})G_9(2x)\\
&+2^{44}x^{992}(1+x^{19}+x^{30}+x^{39})\\
&+2^{44}x^{1095}(1+x^{1100}+x^{2201}).
\end{split}
\]
The rational intermediate in $K_0$ cancels exactly; all four results are positive integers.

\textbf{The 943 indexed fixed tiles.}
Use the actions $u_j,v_j,d_j$ decoded below. For $d=0,1$, let
\[
\begin{gathered}
h(q,b)=2q+b+2,\quad
\mathcal S_d=\{u_j:d_j=d,\ 0\le j<430\},\\
r_{d,q}=205-199d+|\mathcal S_d\cap\rg0{q-1}|.
\end{gathered}
\]
An interval ending at $-1$ is empty. The table defines $t_i$ for
$0\le i\le942$; all index ranges are fixed:
\[
\setlength{\arraycolsep}{1.4pt}\begin{array}{c|l}
\text{index }i&\text{tile }t_i=(w_i,n_i,p_i,s_i)\\\hline
0,1&(2,435,3,1),\ (5,0,0,0)\\
j+2\ (j=2q+b<430)&((1-d_j)r_{0,u_j},h(q,b),d_jr_{1,u_j},v_j)\\
434-2a+b\ (a,b=0,1)&(0,432a+b,0,432a+b)\\
\substack{424+2r_{d,q}+b\\q\in\mathcal S_d,\ d,b=0,1}&(dr_{d,q},b,(1-d)r_{d,q},h(q,b))\\
924+2a+b\ (a,b=0,1)&(0,432a+b,0,438)\\
928&(2,434,2,0)\\
929&(3,436,4,4)\\
930&(4,436,4,0)\\
497+t\ (t=434,435,436)&(1,438,1,t)\\
934+3a+b\ (a,b=0,1,2)&(a^2,437,b,437)
\end{array}
\]
\textbf{A closed decoder for the 430 actions.}
Put $D_b(A,k)=\fl{A/b^k}\bmod b$, where
$a\bmod b=a-b\fl{a/b}$. For $0\le q<215$, $b=0,1$, $j=2q+b$:
\[
\begin{gathered}
N=\mathcal A\bmod2^{636},\quad
T_0=\fl{\mathcal A/2^{636}}\bmod2^{1675},\\
E_0=\fl{\mathcal A/2^{2311}},\\
\nu_j=D_2(N,j),\quad \sigma_j=\sum_{k=0}^{j-1}\nu_k,\\
m_q=D_2(N,430+q),\quad \tau_q=\sum_{k=0}^{q-1}m_k,\\
F_q=14+m_q(D_{16}(E_0,\tau_q)-14),\\
u_j=\nu_j(1+\sigma_j)+(1-\nu_j)D_{216}(T_0,j-\sigma_j),\\
v_j=D_2(F_q,2b),\quad d_j=D_2(F_q,2b+1).
\end{gathered}
\]
Empty sums are zero. Here $(v_0,d_0,u_0)=(1,1,1)$; the constructor tile
uses this actual first target, including its state renaming.

\textbf{The integer $\mathcal A$.}
The alphabet, in order of digit values $0,\ldots,61$, is\par\noindent
{\fontsize{9}{10}\selectfont\ttfamily 0123456789ABCDEFGHIJKLMNOPQRSTUVWXYZabcdefghijklmnopqrstuvwxyz}.\par
Concatenate the following 436 characters as digits $e_0,\ldots,e_{435}$;
put $\mathcal A=\sum_{j=0}^{435}e_j62^{435-j}$.\par
\begingroup\fontsize{9}{10}\selectfont\ttfamily
BxEnhvy6NipzkWmROo0GsA3uK3twWLxscLq20UTumy5Cdu461LwD6g7RDNzF8Y\par
iYDFTpB8ehAypd7d6OsOL4APamnJeNXQypUezCmW9n4oXP6XfgmiWuAWZ8ZJ7U\par
OP7gu4kVdlGZqZHULQMar1dXNeJW4i7oNVI8usp8BlovjjZbwPCovZI328fRUq\par
4q92R81s5GjWQaQfGhkq9UYMK3rRMaBmsOIBkJqc2cj2VRIZ6mjzkBIOkW9SjZ\par
FzeA6J5u4j3sBL3SdgwXzM9mHvQivu8EJ1TBTpRhHDqZIbmYgwQ8VqTAFotl6F\par
We1YC554GKrVHsA8oUGKSA6F6k5KREBSD0Tgd7bFALY84LLbpw574BQyvdoHrl\par
0Oe35uwDf15DpT6VJWZBjSbHPpPAINYTte0eUxnTFTKBlGj8Am6sb1gLldh8WZ\par
8P\par
\endgroup

\endgroup

\end{document}
